\documentclass[11pt]{article}
\usepackage{url}
\usepackage{alltt}
\usepackage{bm}
\usepackage{bbm}
\linespread{1}
\textwidth 6.5in
\oddsidemargin 0.in
\addtolength{\topmargin}{-1in}
\addtolength{\textheight}{2in}

\usepackage{amsmath}
\usepackage{amssymb}
\usepackage{hyperref}

\begin{document}


\begin{center}
\Large
STA 712 Exam 1 Review\\
\normalsize
\vspace{5mm}
\end{center} 

\noindent \textbf{Information:} The midterm exam will cover class material through October 7, and homeworks 1 -- 6. The questions below are not completely comprehensive, but will give you a sense of the kinds of questions I could ask on the exam.

\section*{EDM Questions}

\begin{enumerate}
\item Suppose we are told that $Y \sim EDM(\mu, \phi)$, and we know that $V(\mu) = \mu + \dfrac{\mu^2}{r}$, where $r$ is a \textbf{known} constant. 

\begin{enumerate}
\item Find the canonical parameter $\theta$ as a function of $\mu$.
\item Find the cumulant function $\kappa(\theta)$ as a function of $\theta$.
\end{enumerate}

These values of $\theta$ and $\kappa(\theta)$ match the values for a negative binomial distribution, so if $Y \sim EDM(\mu, \phi)$ and $V(\mu) = \mu + \dfrac{\mu^2}{r}$, then $Y$ must be negative binomial!

\bigskip

\item Determine which of these functions are suitable link functions for a GLM. For those that are not suitable, explain why not.

\begin{enumerate}
\item $g(\mu) = e^\mu$ when $-\infty < \mu < \infty$
\item $g(\mu) = e^{-\mu}$ when $-\infty < \mu < \infty$
\item $g(\mu) = \log \left( \dfrac{\mu}{\mu + 1} \right)$ when $\mu > 0$
\item $g(\mu) = \log(\mu)$ when $-\infty < \mu < \infty$
\end{enumerate}

\bigskip

\item Determine which of the following distributions are EDMs by identifying (where possible) $\theta, \kappa(\theta)$, and $\phi$.

\begin{enumerate}
\item Beta distribution:
$$f(y; a, b) = \frac{\Gamma(a + b)}{\Gamma(a)\Gamma(b)}y^{a-1}(1 - y)^{b-1}$$
for $0 < y < 1$, $a > 0$, and $b > 0$.

\item Strict arcsine distribution:
$$f(y; p) = \frac{A(y)}{y!} \exp\{y \log(p) - \arcsin(p) \}$$
where $y = 0, 1, 2, ...$ and $0 < p < 1$, and $A(y)$ is a complicated function depending only on $y$.
\end{enumerate}

\bigskip

\item
\begin{enumerate}
\item Let $Y \sim Geometric(p)$, with probability function
$$f(y; p) = p(1 - p)^y$$
for $y = 0, 1, 2,...$ and $0 < p < 1$. (This parameterization of the geometric distribution counts the number of failures until the first success, where $p$ is the probability of success). Show that the geometric distribution is an EDM by identifying $\theta$, $\kappa(\theta)$, and $\phi$.

\item Use part (a) and properties of EDMs to find $\mathbb{E}[Y]$ and $Var(Y)$.

\item Find the unit deviance $d(y, \mu)$.

\item Suppose that we observed iid sample $(\mathbf{x}_1, Y_1),...,(\mathbf{x}_n, Y_n)$ from some distribution, where
\begin{align*}
Y_i | \mathbf{x}_i &\sim Geometric(p_i) \\
h(p_i) &= \mathbf{x}_i^T \boldsymbol{\beta}
\end{align*}
If $h(p_i) = g(\mu_i) = \theta_i$, write down either $h(p_i)$ or $g(\mu_i)$.

\item For the model in part (d) (using the canonical link), write down an expression for $U(\boldsymbol{\beta})$ and $\mathbf{H}(\boldsymbol{\beta})$. You do \textbf{not} need to re-derive these from the log-likelihood -- you may rely on general EDM results from class.
\end{enumerate}

\bigskip

\item Recall that the \textbf{cumulant generating function} $C_Y(t)$ for a random variable $Y$ is defined by
$$C_Y(t) = \log \mathbb{E}[e^{tY}]$$
(This definition is not specific to the EDM family). Suppose we have two independent random variables $Y_1$ and $Y_2$ (not necessarily EDMs), for both of which the cgf exists at $t$. Show that
$$C_{Y_1 + Y_2}(t) = C_{Y_1}(t) + C_{Y_2}(t)$$

\bigskip

\item An important property of EDMs (as we have defined them) is that they are \textbf{reproductive}. If $Y_1,...,Y_n \overset{iid}{\sim} EDM(\mu, \phi)$ then
$$\overline{Y} = \frac{1}{n} \sum \limits_{i=1}^n Y_i \sim EDM\left(\mu, \frac{1}{n}\phi\right)$$

\begin{enumerate}
\item In class, we showed that the cgf of $Y \sim EDM(\mu, \phi)$ is 
$$C_Y(t) = \frac{\kappa(\theta + \phi t) - \kappa(\theta)}{\phi}$$
Use this result to show that $\overline{Y} \sim EDM\left(\mu, \frac{1}{n}\phi\right)$.

\item Using part (a), argue that -- for small values of $\phi$ -- if $Y \sim EDM(\mu, \phi)$ then $Y \approx N(\mu, \phi V(\mu))$.
\end{enumerate}

\end{enumerate}

\section*{Logistic regression questions}

\begin{enumerate}
\item[7.] Let $Y_i^*$ be a \textit{continuous} random variable, defined by
$$Y_i^* = \mathbf{x}_i^T \boldsymbol{\beta}^* + \varepsilon_i$$
where the $\varepsilon_i$ are iid samples from the \textbf{Gumbel distribution}, with pdf 
$$f_{\varepsilon}(t) = \exp\{-(t - e^{-t})\} \hspace{0.5cm} t \in (-\infty, \infty)$$
Then, let $Y_i = \mathbb{I}\{Y_i^* > 0\}$. Show that $Y_i | \mathbf{x}_i \sim Bernoulli(p_i)$, with
$$\log(-\log(1 - p_i)) = \mathbf{x}_i^T \boldsymbol{\beta}$$
(where $\boldsymbol{\beta} = - \boldsymbol{\beta}^*$).

The function $\log(-\log(1 - p_i))$ is called the \textbf{complementary log-log} (cloglog) link.

\bigskip

\item[8.] Your friend fits a logistic regression model, and wants to check the quality of the predicted probabilities. They notice that $\frac{1}{n} \sum \limits_{i=1}^n \widehat{p}_i = \frac{1}{n} \sum \limits_{i=1}^n Y_i$ for their fitted model, and they claim this is evidence that the model is well-calibrated. Comment on your friend's interpretation.
\end{enumerate}

\end{document}
